\section{Results}
\label{sec:results}
\label{sec:persistence}

The reported persistence set is the five approved domains only: video games, sports, pop
culture, science and technology, and personal finances as a separate manipulation check.
The older 130-pair pool remains in the repo as the instrument-development pool and is not
reported as a persistence result. The extension run contributes \pqxNEpisodes{} episodes on
\pqxNPairs{} pairs, four arms, three episodes per pair $\times$ arm cell, \pairModel{}.
Every episode reached both elicitations; no refusals, no unparsed choices and no errors in
the reported run.

\subsection{The choice channel: what the challenge asks for decides}
\label{sec:pq1}

\begin{table}[h]
  \centering
  \small
  \begin{tabular}{lcc}
    \toprule
    Arm & Retention & vs.\ control \\
    \midrule
    \texttt{control}               & \pqxeRetCtl{}\%     [\pqxeRetCtlLo{}, \pqxeRetCtlHi{}]         & --- \\
    \texttt{reason\_elicitation}   & \pqxeRetReason{}\%  [\pqxeRetReasonLo{}, \pqxeRetReasonHi{}]   & \pqxeDiffReason{}   [\pqxeDiffReasonLo{}, \pqxeDiffReasonHi{}] \\
    \texttt{self\_critique}        & \pqxeRetCritique{}\%[\pqxeRetCritiqueLo{}, \pqxeRetCritiqueHi{}] & \pqxeDiffCritique{} [\pqxeDiffCritiqueLo{}, \pqxeDiffCritiqueHi{}] \\
    \texttt{counter\_consideration}& \pqxeRetCounter{}\% [\pqxeRetCounterLo{}, \pqxeRetCounterHi{}] & \pqxeDiffCounter{}  [\pqxeDiffCounterLo{}, \pqxeDiffCounterHi{}] \\
    \bottomrule
  \end{tabular}
  \caption{PQ1 on the five approved persistence domains. Retention by arm and the
difference from control, paired within pair. Percentage points, bootstrap 95\% CIs over
\pqxNPairs{} pairs, \num{10000} resamples. Personal finances is labelled separately as a
positive control and not pooled into any PQ estimate.}
  \label{tab:pq1}
\end{table}

Control retains \pqxeRetCtl{}\%. Absent a challenge this procedure returns the same answer
essentially every time, which is what licenses reading the rest as the challenge rather
than instrument noise. Against that baseline \texttt{reason\_elicitation} is
indistinguishable from doing nothing (\pqxeDiffReason{} points,
\([\pqxeDiffReasonLo{}, \pqxeDiffReasonHi{}]\)). \texttt{self\_critique} moves retention by
\pqxeDiffCritique{} points and \texttt{counter\_consideration} by \pqxeDiffCounter{}. The
direction of the request matters. Its mere presence does not. Retention also rises with how
settled the preference already was. The slope of per-pair retention on the extension
balance-pilot consistency is \pqxeSlope{} (\([\pqxeSlopeLo{}, \pqxeSlopeHi{}]\)); that is
PQ2 for the reported persistence set. The remaining domain split and the bias-aware
sensitivity check are in Appendix~\ref{app:persistdetail}.

\subsection{The confidence channel: holding costs confidence}
\label{sec:pq3}

Among episodes where the choice did \emph{not} move, confidence still falls relative to
control: \pqHeldCritique{} points under \texttt{self\_critique}
(\([\pqHeldCritiqueLo{}, \pqHeldCritiqueHi{}]\), \pqHeldCritiquePairs{} pairs) and
\pqHeldCounter{} under \texttt{counter\_consideration}
(\([\pqHeldCounterLo{}, \pqHeldCounterHi{}]\), \pqHeldCounterPairs{} pairs); dropping
episodes that began at zero confidence leaves both unchanged (\pqHeldCritiqueNoZero{},
\pqHeldCounterNoZero{}). This is retention with reduced confidence. The model keeps its
answer and reports itself less sure of it. That is what the design was built to detect, and
what a choice-only measure records as nothing happening. Averaged over \emph{all} episodes
the changes are small and positive (\pqConfReason{}, \pqConfCritique{}, \pqConfCounter{}
points vs.\ control), because held and flipped episodes move in opposite directions. We do
not interpret the flipped half, for the reason given in the Limitations.
Figure~\ref{fig:persistence} in Appendix~\ref{app:persistdetail} shows all three channels.

\subsection{Approved five-domain persistence and the manipulation check}
\label{sec:extension}

The reported persistence set is built from the five approved domains: video games, sports,
pop culture, science and technology, and personal finances as a separate manipulation
check. The original four domains are not reported as persistence results and remain in the
repo as the instrument-development pool only. The extension run therefore uses
\pqxNEpisodes{} episodes on \pqxNPairs{} pairs with the same four-arm protocol and the same
balance-pilot covariate.

The pattern replicates without qualification. On the four preference domains alone,
\texttt{reason\_elicitation} differs from control by \pqxeDiffReason{} points with an
interval of exactly \([\pqxeDiffReasonLo{}, \pqxeDiffReasonHi{}]\). \texttt{self\_critique}
moves retention by \pqxeDiffCritique{} (\([\pqxeDiffCritiqueLo{}, \pqxeDiffCritiqueHi{}]\))
and \texttt{counter\_consideration} by \pqxeDiffCounter{} (\([\pqxeDiffCounterLo{}, \pqxeDiffCounterHi{}]\)). Personal finances is a monotonic arithmetic ladder and is reported separately as the fifth approved domain: it is the manipulation check, not a pooled PQ estimate.

\begin{figure}[t]
  \centering
  \includegraphics[width=\textwidth]{persistence_domains.png}
  \caption{Retention by arm across the four reported preference domains and the separate
  positive-control row: \texttt{video\_games}, \texttt{sports}, \texttt{pop\_culture},
  \texttt{sci\_tech}, and \texttt{finances\_control}. The final row is not pooled into the
  preference-domain estimate; it is shown explicitly as the money-ladder manipulation check,
  where ceiling retention is expected. Bars show the mean per domain and arm with bootstrap
  95\% CIs over pairs.}
  \label{fig:persistence-domains-main}
\end{figure}

\paragraph{PQ4 by domain.}
The domain split reports the four preference domains plus a final row for
\texttt{finances\_control}. The control row is labelled as a positive manipulation check,
never pooled into any PQ estimate. That is the design used in the domain figure: the four
preference domains are the persistence rows, and the fifth row is the money ladder that is
expected to stay at ceiling. We read any movement there as a failure of the elicitation,
not of the preference.

\paragraph{Bias-aware sensitivity check.}
The sports domain is the item type most exposed to the position effect, so the reported
result carries the bias-aware sensitivity check rather than dropping the domain. Excluding
pairs at order gap $\ge \pqxBiasCut{}$ leaves \pqxBiasKeepPairs{} pairs and moves
\texttt{self\_critique} from \pqxeDiffCritique{} to \pqxlDiffCritique{} points and
\texttt{counter\_consideration} from \pqxeDiffCounter{} to \pqxlDiffCounter{}. The
sensitivity check therefore confirms the extension estimate is not being driven by the slot
pattern; it is the challenge effect on the content-bearing items.

\paragraph{Discriminant validity: the control that does not move.}
Personal finances is a monotonic ladder of receive-\$$X$ and owe-\$$X$ outcomes. A forced
choice between two rungs is arithmetic rather than preference, so ceiling retention is the
\emph{expected} result. Wavering there would indicate an unstable elicitation rather than
an unstable preference. It is therefore analysed separately and never pooled into a PQ
estimate. Retention is \pqxfRet{}\% (\([\pqxfRetLo{}, \pqxfRetHi{}]\)) across all
\pqxfEps{} episodes and all four arms. In the collected run the check passes with force.
The two arms that flip \numrange{40}{50}\% of preference pairs move the money ladder not
once: \pqxfFlips{} flips in \pqxfChallengeEps{} challenge episodes. That rules out the
reading that these challenges induce answer-switching regardless of content. The arms move
preferences and leave arithmetic alone. This is the discriminant half of the PQ1 claim, and
the half a preference-only design cannot supply.

\paragraph{Position dependence is item-specific, not only model-specific.}
Section~\ref{sec:validity} established that the order effect varies by model and that
enabling reasoning suppresses it on \pairModel{}. The suppression is not uniform across
items of that same model. In the balance pilot the sports pairs refused nothing, so the
failure mode flagged for them in advance did not occur. They returned a mean per-pair order
gap of \pqxSportsBias{} with reasoning \emph{on}, against \pqxOtherBiasLo{}--%
\pqxOtherBiasHi{} for the other extension domains. Half of them chose the same slot on four
or five of five counterbalanced resamples. Sports is also the domain with the lowest
retention under challenge (\pqxSportsRet{}\% pooled over arms, \pqxSportsRetCritique{}\%
under \texttt{self\_critique}). That is exactly the confound, since an item answered by
slot looks like an item whose preference is easily moved.

No pair was dropped for this. Selecting on a pilot result after seeing it would make the
selection outcome-dependent, so the pilot order gap is carried as a per-pair covariate and
the split is made at analysis time. Excluding the \pqxBiasDropPairs{} extension pairs at an
order gap of \pqxBiasCut{} or above and keeping the remaining \pqxBiasKeepPairs{} moves
\texttt{self\_critique} from \pqxeDiffCritique{} to \pqxlDiffCritique{} points and
\texttt{counter\_consideration} from \pqxeDiffCounter{} to \pqxlDiffCounter{}. Both land
back on the original domains' \pqxoDiffCritique{} and \pqxoDiffCounter{}. Part of the
extension's apparent excess movement was the layout. Once it is removed the two independent
samples agree. The restriction works the other way on $\Delta$confidence, which is why both
are reported. The \texttt{self\_critique} confidence cost is \pqxeConfCritique{} points
across all extension pairs, an interval spanning zero
(\([\pqxeConfCritiqueLo{}, \pqxeConfCritiqueHi{}]\)) and so no measurable difference in this
sample. It is \pqxlConfCritique{} (\([\pqxlConfCritiqueLo{}, \pqxlConfCritiqueHi{}]\)) once
the position-driven items are set aside. We read that conservatively and report it as a
sensitivity check, not a headline. Position-driven items add noise to a between-arm
contrast, so the pooled estimate understates rather than the restricted one overstating.

\subsection{Cross-model replication in the appendix}
\label{sec:crossmodel}

The same arm pattern reproduces on a separate item pool across two further model families:
\texttt{gpt-5.4-nano} and \texttt{gemini-2.5-flash}. The cross-model comparison is
reported in Appendix~\ref{app:persistdetail}; the main text keeps this to a single pointer
sentence because the persistence claim is the five-domain extension on \pairModel{}, and the
arm pattern is checked elsewhere by model family rather than treated as a separate main-text
result.

\subsection{Predictions and outcomes}
\label{sec:predictions}

Three predictions were written into the study design before the persistence run; they were
never tagged as a pre-registration, so Table~\ref{tab:predictions} is a record of what we
expected, not a confirmatory test. The reasoning behind each verdict is in
Appendix~\ref{app:predictions}.

\begin{table}[h]
  \centering
  \small
  \begin{tabular}{p{10.2cm}l}
    \toprule
    Prediction & Verdict \\
    \midrule
    1. \texttt{counter\_consideration} produces the most reversals, \texttt{control} the
    fewest & Half holds \\
    2. \texttt{self\_critique} lowers confidence more than \texttt{reason\_elicitation}
    at equal retention & Not testable as written \\
    3. Higher-consistency pairs retain more & Holds \\
    \bottomrule
  \end{tabular}
  \caption{The three predictions recorded before the persistence run, scored against the
  collected data. Appendix~\ref{app:predictions} gives the reasoning behind each verdict.}
  \label{tab:predictions}
\end{table}
